```latex
\documentclass{article}
\usepackage{amsmath,amssymb,amsthm,amscd}
\usepackage[margin=1in]{geometry}

\newtheorem{theorem}{Theorem}
\newtheorem{lemma}{Lemma}
\newtheorem{proposition}{Proposition}
\newtheorem{corollary}{Corollary}
\theoremstyle{definition}
\newtheorem{definition}{Definition}
\newtheorem{remark}{Remark}

\newcommand{\Uc}{\mathcal{U}}
\newcommand{\bcS}{\mathcal{S}}
\newcommand{\fB}{\mathcal{B}}
\newcommand{\fD}{\mathcal{D}}
\newcommand{\bipr}{\bowtie}
\newcommand{\la}{\langle}
\newcommand{\ra}{\rangle}
\newcommand{\Khat}{\widehat K}
\newcommand{\q}{q}

\begin{document}

\noindent
We work over the ground field $\mathbb C(q)$ with $q$ generic.  
For $1\leq i,j\leq n$, let $e_{ij}$ denote the usual matrix unit.

\section{Preliminaries and the algebra $\bcS_n(\fB^*)$}

Let
\[
R=\q\sum_{i=1}^n e_{ii}\otimes e_{ii}
+\sum_{i\neq j}e_{ii}\otimes e_{jj}
+(\q-\q^{-1})\sum_{i>j}e_{ij}\otimes e_{ji}
\]
be the standard $R$-matrix of type $A_{n-1}$, and let
$R_{21}=P R P$ be its transpose, so that
\[
R_{21}=\q\sum_{i=1}^n e_{ii}\otimes e_{ii}
+\sum_{i\neq j}e_{ii}\otimes e_{jj}
+(\q-\q^{-1})\sum_{i<j}e_{ij}\otimes e_{ji}.
\]
For an $n\times n$ matrix $X=(X_{ij})$ with entries in an associative
algebra, we write as usual
$X_1=X\otimes I_n$ and $X_2=I_n\otimes X$.

\begin{definition}
The algebra $\bcS_n(\fB^*)$ is the associative algebra generated by
the horizontal arcs
\[
b_{ij},\quad c_{ij}\qquad (1\leq i\leq j\leq n),
\]
subject to the following relations. Put
\[
B=(B_{ij})_{1\leq i,j\leq n},\qquad
B_{ij}=
\begin{cases}
b_{ij},& i\leq j,\\
0,& i>j,
\end{cases}
\]
and
\[
C=(C_{ij})_{1\leq i,j\leq n},\qquad
C_{ji}=
\begin{cases}
c_{ij},& i\leq j,\\
0,& i<j.
\end{cases}
\]
Thus $B$ is upper triangular and $C$ is lower triangular. The defining
relations are
\begin{equation}
  R B_1 B_2 = B_2 B_1 R,\qquad
  R C_1 C_2 = C_2 C_1 R,\qquad
  R_{21} C_2 B_1 = B_1 C_2 R_{21}.
  \tag{RL}
\end{equation}
\end{definition}

\begin{remark}
The generators $b_{ij}$ and $c_{ij}$ are exactly the horizontal arcs
directed to the right in the diagram $\bigl(\BigonPGens\bigr)$:
$b_{ij}$ is the upper-right arc from $i$ to $j$, and $c_{ij}$ is the
lower-right arc from $i$ to $j$. The RLL relations (RL) are the
algebraic form of the usual diagrammatic skein relations among such
arcs.
\end{remark}

\section{The Borel subalgebras and their generalized quantum double}

Let $\Uc_q(\mathfrak{sl}(n))=\Uc_q^{\mathrm{ext}}(\mathfrak{sl}(n))$
be the simply-connected quantum group in the convention of
Klimyk--Schm\"udgen. We use the Cartan generators $\Khat_1,\dots,\Khat_n$
with $\Khat_1\cdots \Khat_n$ central, and the simple root vectors
$E_i,F_i$, $1\leq i\leq n-1$. The Borel subalgebras are
\[
\Uc_q(\mathfrak{b}_+)=\bigl\la E_i,\Khat_j^{\pm 1}\mid
1\leq i\leq n-1,\ 1\leq j\leq n\bigr\ra,
\]
\[
\Uc_q(\mathfrak{b}_-)=\bigl\la F_i,\Khat_j^{\pm 1}\mid
1\leq i\leq n-1,\ 1\leq j\leq n\bigr\ra.
\]
Their defining relations are
\[
\Khat_i\Khat_j=\Khat_j\Khat_i,
\]
\[
\Khat_i E_j\Khat_i^{-1}=\q^{\delta_{i,j}-\delta_{i,j+1}}E_j,
\qquad
\Khat_i F_j\Khat_i^{-1}=\q^{-\delta_{i,j}+\delta_{i,j+1}}F_j,
\]
together with the usual $q$-Serre relations among the $E_i$'s and
among the $F_i$'s.

We use the standard Hopf algebra structure on the Borel subalgebras.
Put
\[
K_i=\Khat_i\Khat_{i+1}^{-1}\qquad (1\leq i\leq n-1).
\]
Then
\[
\Delta(E_i)=E_i\otimes 1+K_i\otimes E_i,
\qquad
\Delta(F_i)=F_i\otimes K_i^{-1}+1\otimes F_i,
\]
and the $\Khat_i$'s are group-like.

The Drinfeld--Tanisaki skew pairing between
$\Uc_q(\mathfrak{b}_+)$ and $\Uc_q(\mathfrak{b}_-)$ is given on
generators by
\[
\la \Khat_i,\Khat_j\ra=\q^{\frac1n-\delta_{i,j}},
\qquad
\la E_i,F_j\ra=\frac{\delta_{i,j}}{\q^{-1}-\q},
\]
\[
\la \Khat_i,F_j\ra=\la E_j,\Khat_i\ra=0.
\]
It extends uniquely to a skew Hopf pairing.

\begin{definition}
The generalized quantum double
\[
\fD=\Uc_q(\mathfrak{b}_-)\bipr\Uc_q(\mathfrak{b}_+)
\]
is the algebra generated by the two Borel subalgebras, subject to their
own defining relations and the cross relations
\begin{equation}
  b a=\sum \la a_{(1)},b_{(1)}\ra\, a_{(2)}b_{(2)}\,
  \la S^{-1}(a_{(3)}),b_{(3)}\ra
  \tag{D}
\end{equation}
for $a\in \Uc_q(\mathfrak{b}_+)$ and $b\in\Uc_q(\mathfrak{b}_-)$.
\end{definition}

\section{The $L$-operators and the isomorphism}

We now construct the standard Drinfeld--Reshetikhin--Takahashi
$L$-operators for the two Borel subalgebras.

\begin{lemma}[Standard $L$-operators]\label{lem:L}
There exist unique triangular matrices
\[
L^+=(L^+_{ij})_{1\leq i,j\leq n},\qquad
L^-=(L^-_{ij})_{1\leq i,j\leq n},
\]
with entries in $\Uc_q(\mathfrak{b}_+)$ and $\Uc_q(\mathfrak{b}_-)$
respectively, such that $L^+_{ij}=0$ for $i>j$, $L^-_{ij}=0$ for
$i<j$, and
\[
L^+_{ii}=\Khat_i^+,\qquad
L^+_{i,i+1}=(\q^{-1}-\q)\Khat_i^+ E_i,
\]
\[
L^-_{ii}=(\Khat_i^-)^{-1},\qquad
L^-_{i+1,i}=(\q^{-1}-\q)F_i(\Khat_i^-)^{-1}.
\]
The remaining entries are determined recursively by the RLL relations.
Moreover, the matrices $L^\pm$ satisfy
\begin{equation}
  R L^+_1 L^+_2 = L^+_2 L^+_1 R,\qquad
  R L^-_1 L^-_2 = L^-_2 L^-_1 R,
  \tag{RL-same}
\end{equation}
and the mixed relation
\begin{equation}
  R_{21} L^-_2 L^+_1 = L^+_1 L^-_2 R_{21}.
  \tag{RL-mixed}
\end{equation}
The subalgebra generated by the entries of $L^+$ is
$\Uc_q(\mathfrak{b}_+)$, and the subalgebra generated by the entries
of $L^-$ is $\Uc_q(\mathfrak{b}_-)$.
\end{lemma}

\begin{proof}
This is the standard FRT/RSST construction of the L-operators for the
Borel subalgebras of $\Uc_q^{\mathrm{ext}}(\mathfrak{sl}(n))$; see
Klimyk--Schm\"udgen, {\em Quantum Groups and Their Representations},
Section~9.4. The same-type relations (RL-same) are equivalent to the
defining relations of the Borel subalgebras. The mixed relation
(RL-mixed) is equivalent to the double cross relations (D) for the
given skew pairing. For example, for $n=2$ one obtains from
(RL-mixed) the relations
\[
\Khat_i^+\Khat_j^-=\Khat_j^-\Khat_i^+,
\qquad
E_1F_1-F_1E_1
=\frac{K_{1,+}^{-1}-K_{1,-}}{\q-\q^{-1}},
\]
where $K_{1,\pm}=\Khat_1^\pm(\Khat_2^\pm)^{-1}$. The general case is
obtained by the same computation on adjacent indices and then by the
recursive construction of the remaining entries. 
\end{proof}

\begin{theorem}[Isomorphism with the double]\label{thm:double}
The assignment
\[
\Phi:\bcS_n(\fB^*)\longrightarrow \fD
\]
defined by
\[
\Phi(b_{ij})=L^+_{ij}\quad (1\leq i\leq j\leq n),
\qquad
\Phi(c_{ij})=L^-_{ji}\quad (1\leq i\leq j\leq n)
\]
extends uniquely to an algebra isomorphism.
\end{theorem}

\begin{proof}
By Lemma~\ref{lem:L}, the images $\Phi(b_{ij})$ and $\Phi(c_{ij})$
satisfy the defining relations (RL) of $\bcS_n(\fB^*)$. Hence
$\Phi$ is a well-defined algebra homomorphism.

Conversely, define an algebra homomorphism
\[
\Psi:\fD\longrightarrow \bcS_n(\fB^*)
\]
on the generators of the two Borel subalgebras by
\[
\Psi(\Khat_i^+)=b_{ii},\qquad
\Psi(\Khat_i^-)=c_{ii}^{-1},
\]
\[
\Psi(E_i)=\frac{1}{\q^{-1}-\q} b_{ii}^{-1}b_{i,i+1},
\qquad
\Psi(F_i)=\frac{1}{\q^{-1}-\q} c_{i,i+1}c_{ii}^{-1}.
\]
Again by Lemma~\ref{lem:L}, the relations of the Borel subalgebras
and the double cross relations are satisfied in $\bcS_n(\fB^*)$, so
$\Psi$ is well defined.

It is immediate from the definitions that
\[
\Phi\circ \Psi(E_i)=E_i,\quad
\Phi\circ \Psi(F_i)=F_i,\quad
\Phi\circ \Psi(\Khat_i^\pm)=\Khat_i^\pm,
\]
and that
\[
\Psi\circ \Phi(b_{ij})=b_{ij},\quad
\Psi\circ \Phi(c_{ij})=c_{ij}.
\]
Therefore $\Phi$ and $\Psi$ are mutually inverse algebra isomorphisms.
\end{proof}

\section{The quotient and the quantum group $\Uc_q(\mathfrak{sl}(n))$}

Let $I$ be the two-sided ideal of $\bcS_n(\fB^*)$ generated by
\[
b_{ii}c_{ii}-1\qquad (1\leq i\leq n).
\]

\begin{theorem}\label{thm:quotient}
There is an algebra isomorphism
\[
\bcS_n(\fB^*)/I\ \cong\ \Uc_q(\mathfrak{sl}(n)).
\]
\end{theorem}

\begin{proof}
Under the isomorphism $\Phi$ of Theorem~\ref{thm:double}, the ideal
$I$ is carried to the two-sided ideal
\[
J=\bigl\la L^+_{ii}L^-_{ii}-1\mid 1\leq i\leq n\bigr\ra
=\bigl\la \Khat_i^+(\Khat_i^-)^{-1}-1\mid 1\leq i\leq n\bigr\ra.
\]
Thus
\[
\bcS_n(\fB^*)/I\ \cong\ \fD/J.
\]
In the quotient $\fD/J$, the two families of Cartan generators are
identified:
\[
\Khat_i^+=\Khat_i^-.
\]
After this identification, the same-type RLL relations for $L^+$ and
$L^-$ become the usual RSST presentation of the quantum group
$\Uc_q^{\mathrm{ext}}(\mathfrak{sl}(n))$: the upper and lower
triangular matrices give respectively the upper and lower
$L$-operators of the quantum group, and the mixed RLL relation becomes
the commutator relation between $E_i$ and $F_i$. Up to the standard
automorphism $K_i\mapsto K_i^{-1}$, $E_i\mapsto -E_i$, the latter is
precisely
\[
E_iF_j-F_jE_i
=\delta_{ij}\frac{K_i-K_i^{-1}}{\q-\q^{-1}},
\qquad
K_i=\Khat_i\Khat_{i+1}^{-1}.
\]
The $q$-Serre relations are already contained in the same-type RLL
relations. Therefore the quotient is exactly the simply-connected
quantum group $\Uc_q(\mathfrak{sl}(n))$, as presented in
Klimyk--Schm\"udgen, Section~9.4.
\end{proof}

\begin{remark}
The diagonal relations $b_{ii}c_{ii}=1$ are precisely the relations
that identify the two Cartan copies in the generalized double. This
is why the double of the two Borel subalgebras becomes the quantum
group itself after imposing these relations.
\end{remark}

\end{document}
```